% !TeX root = ../../Book3.tex
% 纲要标题：### 6.5 正规化有限性的证明
% 纲要位置：计划纲要.md，第 1671 行。

\section{正规化有限性的证明}
\label{b3:sec:0605}

% TODO: 按以下纲要要点撰写本节。

% 纲要内容（仅作写作计划，尚非教材正文）：
%
% * 用 Noether 正规化把有限型整环的整闭包问题归约为多项式环在有限分式域扩张中的整闭包问题
% * 对有限可分扩张建立迹配对和公共分母控制引理，利用多项式环的整闭性与 Noether 性得到有限生成
% * 正特征下单独建立不可分情形所需的有限性引理；通过有限个系数和纯不可分扩张的代数处理完成证明，不把可分迹论证无条件套用，也不默设基域完美
% * 从整环推广到约化有限型代数的有限个极小素分支；明确总分式环中的正规化与各分支整闭包的关系
% * 本节接续第6.4节的正规化对象，准确陈述并完整证明有限性定理；不以附录中的一般 Nagata 环定理代替所需引理。结论范围核对见 [V13-2]
%

\subsection{由 Noether 正规化归约整闭包问题}
\label{b3:subsec:0605:01}

待撰写。

\subsection{可分扩张中的迹配对与公共分母}
\label{b3:subsec:0605:02}

待撰写。

\subsection{正特征不可分扩张中的有限性引理}
\label{b3:subsec:0605:03}

待撰写。

\subsection{约化代数的极小素分支与总分式环}
\label{b3:subsec:0605:04}

待撰写。

\subsection{正规化有限性定理的完整证明}
\label{b3:subsec:0605:05}

待撰写。
